% Số: facts.md mục Test-4 ("Hạn chế phải nêu (v3)": lực 26%, ~120 họ, đọc trái phép v1 20 / v3 18 ca, tiện ích tuyệt
% đối thấp; "40 ca vượt v1" trong mẫu 38 -> 2; M3 Test-II 9 cặp so với Test-4 1), mục "Bổ sung cho bài v3" (phát lại
% oracle Dev-val), mục baseline bổ sung (Prompt Guard 0 lần chặn), mục AgentDojo v3.
\section{Thảo luận}\label{sec:discussion}

\paragraph{Vì sao giả thuyết chính không đạt.} H1 được thiết kế với giả định \base{} để lọt khoảng 14 trong 48 họ,
như trên \tTwo. Trên \tThree, \base{} chỉ để lọt 6 họ. Khác biệt nằm chủ yếu ở rò rỉ qua sink: trên \tTwo{} cơ chế này
có 9 cặp họ và mô hình bị hại dưới \base, trên \tThree{} chỉ còn 1. Với nền thấp như vậy, ngay cả khi \ours{} chặn được
mọi họ mà \base{} để lọt, số cặp lệch vẫn rất nhỏ. Lực của kiểm định ở mức nền này chỉ khoảng 26\%, và cần khoảng 120
họ để đạt lực 80\%. Kết quả trên \tThree{} vì vậy là một hiệu ứng sàn chứ không phải bằng chứng rằng hai cấu hình như
nhau. Các nguồn khác cho cùng chiều. Trên \dev, số cặp họ và mô hình bị hại giảm từ 36 xuống 4. Chấm lại 40 ca từng lọt
qua \base{} trên hai tập niêm phong trước cho 38 ca bị hại dưới \base{} và 2 dưới \ours, cả hai đều là sửa tham số.
Tuy vậy, các số này đến từ dữ liệu đã dùng khi phát triển \ours, nên chỉ có giá trị chẩn đoán và không thay được kiểm
định đăng ký trước. Bài báo không thêm mô hình hay chạy lại để tìm ý nghĩa thống kê sau khi đã biết kết quả.

\paragraph{Mỗi thành phần đóng góp gì.} Bằng chứng mạnh nhất cho \cOne{} và \cTwo{} là bằng chứng mức ca. Bốn họ mà
\base{} để lọt đều bị cổng của \ours{} chặn trước khi thực thi, và mỗi họ lọt trở lại khi bỏ đúng thành phần đã chặn
nó. \cTwo{} xử lý đúng chỗ hở lớn nhất của \base, là giá trị mật nằm trong nguồn gắn nhãn công khai. Dưới \ours, không
giá trị nào lộ ở câu trả lời cuối. Tác dụng của \cOne{} khi chạy thật nhỏ hơn khi chấm offline. Trên các lời gọi sink
của quỹ đạo tham chiếu thuộc phần kiểm định của \dev, \cOne{} chặn đủ 35 lời gọi tác dụng phụ ngoài yêu cầu, so với 15
của \rDest. Nhưng các mô hình hiếm khi tự đề xuất những đích này, nên trên \tThree{} \cOne{} chỉ tạo khác biệt ở 2 họ.
\cThree{} chưa có tác dụng đo được trên \tThree. Hai họ còn lọt là sửa giá trị số hoặc đơn vị. Có hai cách giải thích, đều chưa được kiểm từng ca.
Hoặc giá trị bị sửa không nằm trong tập literal được ghim, tức (A5) bị vi phạm. Hoặc literal vẫn có mặt cạnh một giá
trị khác, trường hợp Mệnh đề~\ref{prop:int} không loại trừ. Cả hai đều cho thấy \cThree{} cần ngữ nghĩa của trường
thay vì chỉ kiểm sự có mặt của literal.

\paragraph{Chi phí tiện ích.} \ours{} không làm tiện ích thấp hơn \base{} quá 0,6 điểm ở nhánh nào, nhưng cả hai đều
thấp hơn \cNone{} khoảng 3 điểm. Một phần nguyên nhân là chặn nhầm. \ours{} chặn nhầm 24 trong 960 ca lành tính của
\tThree, so với 19 của \base. Trên \dev, lý do chặn nhầm thường gặp là các luật giữ từ \base{} (\rFlow{} theo nguồn và
\rSig). Riêng với Llama-3.1, \cThree{} đứng đầu, vì mô hình này hay viết lại giá trị người dùng đã nêu khi soạn nội
dung gửi. Tiện ích tuyệt đối thấp ở mọi cấu hình, kể cả \cNone{} (42,8\%), do giới hạn của các mô hình nhỏ như hết bước
hay đầu ra bị cắt, không do phòng thủ. Trên AgentDojo, chi phí tập trung ở các nhiệm vụ lấy đích từ dữ liệu
công cụ. Trong benchmark, ``cần xác nhận'' được thi hành như từ chối. Một triển khai thật có thể hỏi người dùng ở bước
này, và đây là cách hợp lý để lấy lại tiện ích mà không cấp quyền cho nội dung không đáng tin.

\paragraph{Ngôn ngữ.} Luật \rDest{} của \base{} so khớp chuỗi nên không phụ thuộc ngôn ngữ. \cOne{} thì có, vì cực tính
dựa trên danh sách từ phủ định, loại trừ và tương phản cho tiếng Việt có dấu, không dấu và tiếng Anh. Một ngôn ngữ hay
phương ngữ mới cần bổ sung danh sách này, và một phủ định ngoài danh sách sẽ làm cổng cho qua một đích lẽ ra bị cấm.
Thiết kế giảm rủi ro này bằng cách trả về \emph{cần xác nhận} khi phủ định không rõ phạm vi, nhưng không loại bỏ được
nó. Với phòng thủ phát hiện, vấn đề ngôn ngữ còn rõ hơn: \cPG{} không gắn cờ một quan sát nào của \tThree, tức bộ phân
loại này không chuyển được sang dữ liệu tiếng Việt của benchmark.

\paragraph{Mối đe dọa với tính hợp lệ.} Về cấu trúc đo, vị từ gây hại và bộ chấm được đóng băng trước khi chạy, nhưng
dữ liệu là tổng hợp, công cụ là bản giả và mỗi nhiệm vụ chỉ có một lượt. Về tính hợp lệ nội tại, \tThree{} được sinh
từ đặc tả bằng mã và được chính nhóm phát triển phòng thủ rà soát. Tập được niêm phong trước khi có mã của \ours, nhưng
không loại trừ được việc người soạn vô tình chọn các ca hợp với thiết kế. Bộ thích nghi thứ hai do thành viên của nhóm
soạn, chỉ có một vòng và chạy trên tập phát triển. Năm baseline bổ sung được chạy sau khi đã biết kết quả chính. Về
tính hợp lệ bên ngoài, các mô hình mở chỉ có 8 đến 12 tỉ tham số, Gemini 3 Flash chạy qua proxy và không ghim được
phiên bản, còn AgentDojo chỉ chạy với một mô hình.

\paragraph{Hạn chế.} Các mệnh đề được chứng minh trên đặc tả, chưa được kiểm hình thức trên mã. Kiểm thử dựa trên thuộc
tính với một văn phạm sinh yêu cầu và chuỗi lời gọi là bước tiếp theo tự nhiên. Đọc trái phép nằm ngoài phạm vi: dưới
\ours{} vẫn có 18 trong 720 ca tấn công có lời gọi đọc trái phép, gần như toàn bộ trên Llama-3.1 (\base: 20). \cTwo{}
không theo dõi giá trị qua các phép biến đổi ngoài (A4), như mã hóa hay chia nhỏ qua nhiều lời gọi. Bộ chấm biến đổi
không tìm thấy ca nào như vậy trên dữ liệu hiện có, nên giả định này chưa thực sự bị thử thách. Cuối cùng, bộ dò của
\cTwo{} được thiết kế để không phụ thuộc vào dạng giá trị mật của bộ sinh dữ liệu, nhưng độ phủ trên giá trị mật thật
vẫn phải đo trên dữ liệu thật.
